\documentclass{article}
\usepackage{amsmath}
\usepackage{amssymb}
\usepackage{amsthm}
\usepackage[T1]{fontenc}

\newtheorem*{conjecture}{Conjecture}
\newtheorem*{theorem}{Theorem}
\newtheorem{lemma}{Lemma}
\newtheorem*{lemma*}{Lemma}
\newtheorem*{proposition}{Proposition}


\makeatletter
\newcommand*{\rom}[1]{\expandafter\@slowromancap\romannumeral #1@}
\makeatother


\begin{document}


\section{Zadanie}

Udowodnij, że dla każdej liczby rzeczywistej \(x \in (16, 81)\) zachodzi nierówność:

\begin{equation*}
    \sqrt[4]{x} > \frac{x + 114}{65}.
\end{equation*}

\section*{Rozwiązanie}

Chcemy pokazać, że funkcja \(f(x) = \sqrt[4]{x} - \frac{x + 114}{65}\) jest dodatnia na przedziale \((16, 81)\).
Posłużymy się, związkami pochodnej z monotonicznością funkcji. Po krótkich rachunkach, otrzymujemy:

\begin{equation*}
    f'(x) = \frac{1}{4 \sqrt[4]{x^{3}}} - \frac{1}{65}
\end{equation*}

Pochodna \(f'(x)\) jest malejąca na przedziale (16, 81). Istotnie, dla wszystkich \(a, b \in (16, 81)\) takich, że
\(a < b\) mamy:

\begin{equation*}
    a < b
\end{equation*}
\begin{equation*}
    4 \sqrt[4]{a^{3}} < 4 \sqrt[4]{b^{3}} 
\end{equation*}
\begin{equation*}
    \frac{1}{4 \sqrt[4]{a^{3}}} > \frac{1}{4 \sqrt[4]{b^{3}}}
\end{equation*}
\begin{equation*}
    f'(a) = \frac{1}{4 \sqrt[4]{a^{3}}} - \frac{1}{65}  > \frac{1}{4 \sqrt[4]{b^{3}}} - \frac{1}{65} = f'(b)
\end{equation*}

Następnie rozpatrujemy wartości pochodnej na przedziale \((16, 81)\). Niech \(x_{0}\) będzie miejscem zerowym
pochodnej \(f'(x)\). Dla każdego \(c \in (16, x_{0})\) mamy:

\begin{equation*}
    f'(c) > f'(x_{0}) = 0
\end{equation*}
a dla każdego \(d \in (x_{0}, 81)\) mamy:
\begin{equation*}
    0 = f'(x_{0}) > f'(d)
\end{equation*}
zatem \(f\) jest rosnąca na \([16, x_{0}]\) oraz malejąca na \([x_{0}, 81]\). Dla \(x \in (16, x_{0}]\) mamy
\(0 = f(16) < f(x)\). Dla \(x \in [x_{0}, 81)\) mamy \(f(x) > f(81) = 0\). Ostatecznie:

\begin{equation*}
    f(x) > 0 \ \mbox{dla} \ x \in (16, 81)
\end{equation*}

\end{document}